
% ============================================================================
\clearpage
\section*{v4.10.0-r65 Source-Only: Charge Axis and Charge Type on the Sedenion Carrier; Route A and the Relative CP Sign Depend on One Named Setting}
\addcontentsline{toc}{section}{v4.10.0-r65 Source-Only: Charge Axis and Charge Type on the Sedenion Carrier; Route A and the Relative CP Sign Depend on One Named Setting}
\label{sec:v4100-r65-charge-axis-type}
% ============================================================================

\paragraph{Purpose.}
Layer r64.2 recorded that H130 does not support the charge axis \(e_8\) (placeholder CT-003) and left open whether route A of r13 (two generations)
and the relative CP sign of r59.6 survive on the axis \(e_7\).  This layer records the test H162 (k360, 26 November 2026), which answers that
question, and the exploration k361, which adds the second type of charge and finds that the relative sign of r59.6 rests on a setting that was
not named there.  It restricts one sentence of r59.6 and weakens the condition under which r13.3, r13.4 and r59.6 stand.  Earlier layers are not
rewritten.  No guard changes.  No comparison with data is made.

\noindent\textbf{Source-only status.}
This v4.10.0-r65 layer is source-only.  It does not build a PDF, does not assign a new DOI,
and does not authorize a public release.  The v5.4.0 Version DOI
\href{https://doi.org/10.5281/zenodo.24500009}{10.5281/zenodo.24500009} and the Concept DOI
\href{https://doi.org/10.5281/zenodo.19210728}{10.5281/zenodo.19210728} are unchanged by it.

\subsection*{v4.10.0-r65 results}
\label{subsec:v4100-r65-results}

Notation.  \(H_n=\operatorname{span}(1,e_7,e_8,e_{15})\) is the colour-neutral block of \(\mathbb S\); it is associative (a quaternion algebra).
\(P_0\), \(P_3\), \(P_{\bar3}\) are the projectors onto \(H_n\) and the two coloured blocks (ranks \(4,6,6\)).  For a unit
\(u\in\operatorname{Im}H_n=\operatorname{span}(e_7,e_8,e_{15})\),
\[
Q^{L}_u=P_0\bigl(\tfrac12-\tfrac i2L_u\bigr)P_0+\tfrac23P_3+\tfrac13P_{\bar3},
\qquad
Q^{R}_u=P_0\bigl(\tfrac12-\tfrac i2R_u\bigr)P_0+\tfrac23P_3+\tfrac13P_{\bar3},
\]
with left and right multiplication \(L_u\), \(R_u\).  The charge of r13.1 is \(Q_*=Q^L_{e_8}\).

\paragraph{Theoremlet v4.10.0-r65.1 (the charge has the same spectrum and the same symmetries for every axis; H162a, not blind).}
For every unit \(u\in\operatorname{Im}H_n\), \(Q^L_u\) has the spectrum \(0\,(2)\), \(\tfrac13\,(6)\), \(\tfrac23\,(6)\), \(1\,(2)\), commutes with
\(su(3)_c\) and satisfies \(\overline{Q}=1-Q\).  \textbf{Confirmed.}
\emph{Proof sketch.}  \(H_n\) is associative, so \(L_u^2=-1\) on \(H_n\) for every unit \(u\), and \(\tfrac12-\tfrac i2L_u\) has the eigenvalues
\(0\) and \(1\) twice each.  The coloured part of \(Q^L_u\) does not depend on \(u\).
\emph{Check.} Wolfram Language exact for \(u=e_7\): characteristic polynomial \(x^2(x-1)^2(3x-1)^6(3x-2)^6/3^{12}\), both commutation residues
exactly zero; \(L_u^2=-1\) on \(H_n\) symbolically in \(u\).  numpy: 200 of 200 random axes.

\paragraph{Theoremlet v4.10.0-r65.2 (route A, the family \(su(2)\) and stage A of H125 hold for every axis; H162b, c; b not blind).}
For every unit \(u\in\operatorname{Im}H_n\): the right multiplications \(R_x\), \(x\in\{e_7,e_8,e_{15}\}\), commute with \(Q^L_u\) and with
\(su(3)_c\); \([R_{e_7},R_{e_8}]=-2R_{e_{15}}\) on the lepton block and \(+2R_{e_{15}}\) on the coloured block; the Bargmann invariant of the octant
projectors is \(\tfrac14-\tfrac i4\) on the leptons against \(\tfrac14+\tfrac i4\) on the quarks; every particle species is a doublet (Casimir
\(\tfrac34\)), so the carrier holds two complete generations.  With \(Q^L_{e_7}\) the space of intertwiners to
\(\mathbb C\otimes\mathbb H\otimes\mathbb O\) has dimension 20, and all 24 assignments of full rank per sector have determinant \(-1\) (leptons) and
\(+1\) (quarks), as for \(e_8\) (k252); the same for the two rotated copy axes.  H162b and H162c \textbf{confirmed}.
\emph{Reason.}  Right multiplications commute with \(L_u\) on \(H_n\) by associativity, and with \(P_3\) because they commute with \(su(3)_c\).  The
family algebra and its opposite orientation on the two sectors are properties of the sedenion right multiplication and of colour.  The charge
axis does not enter.
\emph{Check.} Commutators and Bargmann invariants Wolfram Language exact, built independently from the Cayley--Dickson recursion; part c numerically
only (numpy, residues below \(10^{-10}\)).

\paragraph{Theoremlet v4.10.0-r65.3 (what distinguishes \(e_8\) and what distinguishes \(e_7\); H162d, not blind, and exploration H162e).}
\begin{enumerate}[leftmargin=2em]
\item[(a)] \(Q^L_{e_7}\) leaves \(\mathbb C\otimes\mathbb O_0\) invariant (leak \(0\)) and equals there the octonionic charge \(Q_8\) of one complete
generation (deviation \(0\), exact).  \(Q_*\) does not: squared leak \(\tfrac12\).  On the other two copies \(\mathbb O_{120}\), \(\mathbb O_{240}\) the
squared leak of \(Q^L_{e_7}\) is \(\tfrac38\).  Each of the three copy axes does the same for its own copy.  H162d \textbf{confirmed}.
\item[(b)] The \(S_3\) cycle \(\psi\) commutes with \(Q^L_u\) only for \(u=\pm e_8\); \(\lVert[Q^L_{e_7},\psi]\rVert=\sqrt3\), and \(\psi\) does not
map the neutral kernel of \(Q^L_{e_7}\) into itself.
\item[(c)] No axis satisfies both.  The statements that use \(\psi\) on the neutral kernel (r13.1, r13.2, the phase per generation in r13.3) do not
hold in that form for \(e_7\).  A substitute cycle \(\exp(-(\pi/3)R_u)\) reproduces \(\rho^2=\tfrac12\) and the phase \(\pi/4\) for \(u=e_7\) and
\(u=e_8\) (for the rotated copy axes \(\rho^2=0.283\) or \(0.717\)); it is a set right rotation and not an automorphism.  As in r13.2 this is a
construction and not a derivation of Koide.
\end{enumerate}
The choice of the axis stays a setting, now between two named criteria: \(\psi\)-fixed (\(e_8\)) or faithful to a copy (\(e_7\) and its two
rotations).

\paragraph{Theoremlet v4.10.0-r65.4 (two types of charge, and the side that carries the family; exploration k361, not blind, numerical).}
\begin{enumerate}[leftmargin=2em]
\item[(a)] \emph{Completeness.}  A Hermitian charge on \(\mathbb C\otimes H_n\) with spectrum \(\{0,0,1,1\}\) and \(\overline Q=1-Q\) has the form
\(\tfrac12-\tfrac i2A\) with a real antisymmetric \(A\), \(A^2=-1\).  On \(H_n\), \(so(4)=\{L_a+R_b\}\), and \(A^2=-1\) holds exactly for
\(A=L_u\) or \(A=R_u\) with a unit \(u\).  There are two types, each with a free axis (r13.1 lists the second as ``up to two signs'').
\item[(b)] \(Q^R_u\) satisfies the same conditions as \(Q^L_u\): the spectrum of r65.1, \([Q,su(3)_c]=0\), \(\overline Q=1-Q\); \(Q^R_{e_8}\)
commutes with \(\psi\), and \(Q^R_{e_7}\) extends \(Q_8\) on \(\mathbb O_0\) (leak \(0\)).  No criterion of H162 separates the types.
\item[(c)] With \(Q^R_u\) the \emph{left} multiplications \(L_x\) commute with the charge on the leptons and the right multiplications do not
(commutator norm \(2\)).  On the coloured blocks both \(L_x\) and \(R_x\) commute with charge and colour.
\item[(d)] Orientation, as the sign \(s\) in \([X_{e_7},X_{e_8}]=2sX_{e_{15}}\) with the Bargmann invariant \(\tfrac14+\tfrac{is}4\):
\begin{center}
\small
\begin{tabular}{@{}lll@{}}
\toprule
block & \(X=L\) & \(X=R\)\\
\midrule
leptons (\(H_n\)) & \(s=+1\) & \(s=-1\)\\
quarks (coloured) & \(s=-1\) & \(s=+1\)\\
\bottomrule
\end{tabular}
\end{center}
On the quarks there are two native family \(su(2)\) of opposite orientation.  Which one is admissible on the leptons is fixed by the type of charge.
\item[(e)] If \emph{one} native multiplication carries the family on the whole carrier (\(R_x\) for the left type as in r13.4, \(L_x\) for the right
type), leptons and quarks are oppositely oriented, for both types and every axis.  If the two sectors use different sides, which charge and colour
permit because both sides commute on the quarks, they are equally oriented: computed for \(Q^R\) with \(L\) on the leptons and \(R\) on the
quarks, \(24+24\) assignments, determinants \(+1\) and \(+1\).
\end{enumerate}
\emph{Check.} numpy on the infrastructure of r13, r14 and k252 (residues below \(10^{-10}\)); part (a) is the standard description of the orthogonal
complex structures of \(\mathbb R^4\).  No Wolfram cross-check was made for k361.

\subsection*{v4.10.0-r65 precisions to earlier layers}
\label{subsec:v4100-r65-precisions}

\begin{enumerate}[leftmargin=2em]
\item[(a)] r64.2: the question ``whether they survive on \(u=e_7\) is not computed'' is answered by r65.2: they do.  H130 does not endanger route A,
the family \(su(2)\) or stage A of H125.
\item[(b)] r13.3, r13.4, r14.4, r59.6: the condition ``under CT-003'' can be weakened to ``under a charge of the form \(Q^L_u\) or \(Q^R_u\) and the
setting S-side below''.  The statements of r13.1--r13.3 that use \(\psi\) stay conditional on the axis \(e_8\) (r65.3).
\item[(c)] r59.6: the sentence that the relative sign ``is the first item of the series that does not set a sign'' is restricted.  The relative sign
does not depend on the convention \(i\mapsto-i\), on the charge axis or on the charge type.  It does depend on the setting
\textbf{S-side: the same side of the multiplication carries the family in both sectors.}  S-side is the natural choice (r13.4: one native operation), but
it does not follow from charge and colour alone (r65.4e).
\item[(d)] Stage B of H125 (\(\sin\delta_{\rm PMNS}<0\), deadline 31 December 2035) now depends on S1 and S-side.  Its status ``open, prior probability
\(\tfrac12\)'' and its decision rule are unchanged.
\end{enumerate}
Not computed: whether an action (W1/W2, Yukawa connectors) fixes the side; whether r14.4 keeps all its statements beyond r65.2 on the axis \(e_7\).

\subsection*{v4.10.0-r65 frames}
\label{subsec:v4100-r65-frames}

\begin{center}
\small
\begin{tabular}{@{}>{\raggedright\arraybackslash}p{0.34\textwidth}>{\raggedright\arraybackslash}p{0.20\textwidth}>{\raggedright\arraybackslash}p{0.38\textwidth}@{}}
\toprule
item & frame & comparison\\
\midrule
r65.1--r65.4 (charge, family algebra, orientation) & frame independent algebra on the carrier & none with data\\
complex conjugation \(\overline Q=1-Q\) & charge conjugation of the carrier & not identified with \(z\mapsto\bar z\) or \(z\mapsto iz\) of the frames\\
stage B of H125 & M (measured mixing phases) & open; only after the declared map of r64.3\\
T, P & not used & ---\\
\bottomrule
\end{tabular}
\end{center}

\subsection*{v4.10.0-r65 register entries}
\label{subsec:v4100-r65-register}

\begin{center}
\small
\begin{tabular}{@{}>{\raggedright\arraybackslash}p{0.40\textwidth}>{\raggedright\arraybackslash}p{0.52\textwidth}@{}}
\toprule
statement & status after v4.10.0-r65\\
\midrule
H162a, b, d & confirmed, not blind\\
H162c & confirmed (numerical)\\
H162e, k361 & explorations\\
route A (r13.3), family \(su(2)\) (r13.4), H125 stage A (r59.6) & hold for every axis and both charge types; conditional on S-side\\
H125 stage B & open, deadline 31 December 2035; depends on S1 and S-side\\
CT-003 (charge axis) & placeholder; choice between \(\psi\)-fixed (\(e_8\)) and faithful to a copy (\(e_7\)) not decided\\
S-side (same side in both sectors) & setting, named here; not derived\\
number of generations & two on this carrier, not three\\
OP1--OP5 & open\\
\bottomrule
\end{tabular}
\end{center}

\subsection*{v4.10.0-r65 notes}
\label{subsec:v4100-r65-notes}

{\raggedright
Advance and result note of H162: \texttt{notes/\allowbreak research/\allowbreak predictions/\allowbreak Vorab\_\allowbreak Notiz\_\allowbreak H162\_\allowbreak Ladungsachse\_\allowbreak e7\_\allowbreak WegA\_\allowbreak CP.md}
(commit \texttt{080d9a9}, before the computation) and \texttt{Ergebnis\_\allowbreak H162\_\allowbreak Ladungsachse\_\allowbreak e7\_\allowbreak WegA\_\allowbreak CP.md}.
Scripts and actual outputs: \texttt{notes/\allowbreak research/\allowbreak hilbert\_\allowbreak shift\_\allowbreak felder\_\allowbreak 2026\_\allowbreak 11\_\allowbreak 24/\allowbreak ladungsachse\_\allowbreak h162/}
(\texttt{k360\_\allowbreak test\_\allowbreak h162.py}, \texttt{k360b\_\allowbreak exakt.wl}, \texttt{k361\_\allowbreak rechtsladung.py} with their output files and \texttt{README.md}).
Disclosed there: the first Wolfram run of k360 was invalid (inverse of a singular matrix, a clash of symbol names) and was repeated.
Tools: Python (numpy) and Wolfram Language; Qiskit not used (no circuit); no HDL comparison (no finite search or guard procedure involved).
\par}

\subsection*{v4.10.0-r65 ledger and guard update}
\label{subsec:v4100-r65-ledger-guard}

\[
R_{\rm chk}^{v4100\text{-}r65}=(1,1,0)
\]
for (status recorded; computations checked in the cited scripts; no comparison with data).

The layer spends no guard:
\[
R_{\rm glob}^{(v4100\text{-}r65,+)}=(0,1,1,1,1),
\qquad
\Gamma_{\rm guard}^{v4100\text{-}r65}=(1,1,1,1,1),
\qquad
\beta_4=1/256\ \text{unspent}.
\]
The layer does not imply exact finite \(C5\) zero closure, exact all-readout terminality, \(C6/T2\), a global photon-clock/RH theorem,
or global all-readout closure.  OP1--OP5, C5 and C6/T2 remain open.
